% ============================================================
%  ABSTRACT
% ============================================================
\begin{abstract}
Los ciclones extratropicales (ETC) generan vientos intensos y precipitaciones cuantiosas en la costa del Atlántico Sudoccidental, cuya coincidencia espacial suele asumirse sin verificación explícita. Mientras el Hemisferio Norte cuenta con estudios sistemáticos sobre cómo se distribuyen estos campos, el Hemisferio Sur carece de una integración equivalente que combine densidad espacial, análisis de Funciones Ortogonales Empíricas (EOF) univariadas y multivariadas (MEOF) por fase del ciclo de vida. Esta tesis desarrolla dicha integración para evaluar si la intensidad dinámica del ETC predice la magnitud y ubicación de sus máximos de viento a 10 metros y precipitación, o si ambos campos desarrollan una organización estructural cualitativamente distinta según la fase evolutiva.

El análisis se construye sobre el catálogo de trayectorias por vorticidad relativa a 850~hPa \citep{gramcianinov2020analysis} y la clasificación objetiva del ciclo de vida en cuatro fases mediante \textit{CycloPhaser} \citep{coutodesouza2024new, deSouza2025CycloPhaser}, sobre ERA5 (1979--2020). Se estratificó la muestra por tres regiones de ciclogénesis (SBR, LPB, ARG) y dos estratos de intensidad: la Población Global y el subconjunto de ciclones más intensos (p90, extraído de la Población Global). Sobre esta base se aplicaron, de forma progresiva, funciones de densidad de probabilidad (PDF) de magnitud, estimación de densidad espacial (PDFe) de ubicación, descomposición EOF univariada con evaluación de significancia mediante Bootstrap-$t$, y descomposición MEOF de la covarianza conjunta, integradas finalmente en prototipos estructurales acoplados.

En la Población Global, el viento se comporta de forma pasiva y predecible, con un único modo EOF dominante (31--57\% de la varianza) acoplado a la vorticidad; la precipitación, en cambio, se distribuye entre múltiples modos (EOF1 apenas 11--16\%). Esta arquitectura se reorganiza drásticamente en el subconjunto p90: el viento se hiperconcentra en el cuadrante noroeste (moda de hasta 23~m\,s$^{-1}$, densidades de hasta $27\times10^{-4}$ en SBR), mientras la precipitación colapsa en el núcleo y se reorganiza en una banda enroscada hacia el sureste y sur. El MEOF muestra que ambas transformaciones son una misma reorganización: en la Población Global captura 19--27\% de la covarianza conjunta con cargas co-localizadas, pero en p90 cae a 12--19\% y revela un patrón de anticovarianza geométrica, respaldado por correlaciones negativas entre las componentes principales de ambos campos ($r$ entre $-0.31$ y $-0.45$). Los prototipos estructurales cuantifican esta separación en más de $8^{\circ}$--$9^{\circ}$ durante la madurez y el decaimiento, con ARG alcanzando la mayor magnitud (${\sim}10$--$11^{\circ}$) pese a su oclusión más gradual. SBR y LPB, sostenidos por flujos de calor latente oceánico y humedad del SALLJ, desarrollan el desacople más rápido; ARG, dominado por dinámica baroclínica y modulación orográfica, lo hace de forma más gradual pero no menos pronunciada.

Estos resultados muestran que, en los ETC maduros del Atlántico Sudoccidental, los máximos de viento y precipitación operan en cuadrantes opuestos, y que su parametrización mediante métricas escalares de intensidad central subestima sistemáticamente tanto la magnitud como la ubicación de ambos forzantes. Hasta donde se ha podido verificar, este trabajo constituye la primera caracterización que combina PDFe espacial, EOF univariado y MEOF para cuantificar este desacople por fase evolutiva y región de génesis en el Hemisferio Sur, con implicaciones para la gestión del riesgo costero en Brasil, Uruguay y Argentina.
\end{abstract}

\vspace{1cm}

\noindent \textbf{Keywords:} South America; extratropical cyclones; wind at 10 meters; precipitation; spatial structure; multivariate EOF.
